\documentclass[runningheads]{llncs}
\usepackage{graphicx}
\usepackage{booktabs}
\usepackage{hyperref}
\usepackage[utf8]{inputenc}
\usepackage[T1,T2A]{fontenc}
\usepackage[english]{babel}
\usepackage{array}
\usepackage{float}

\newcolumntype{P}[1]{>{\raggedright\arraybackslash}p{#1}}

% Optimize float placement parameters to keep tables tightly placed with text
\renewcommand{\topfraction}{0.90}
\renewcommand{\bottomfraction}{0.85}
\renewcommand{\textfraction}{0.05}
\renewcommand{\floatpagefraction}{0.85}

\begin{document}
\sloppy

\title{Detecting AI-Generated User Reviews in Kazakh: A Study on BERT Model Performance and False Positive Reduction}
\titlerunning{Detecting AI-Generated Reviews in Kazakh}

\author{Daulet Anekesh\inst{1} \and Irina Ualiyeva\inst{1}}
\authorrunning{D. Anekesh and I. Ualiyeva}

\institute{Al-Farabi Kazakh National University, Almaty, Kazakhstan \\
\email{\{anekeshd, i.ualiyeva\}@gmail.com}}

\maketitle

\begin{abstract}
The rapid proliferation of Large Language Models (LLMs) has intensified the need for reliable machine-generated text detection. However, research remains sparse for morphologically rich and agglutinative languages such as Kazakh. In this paper, we present an empirical benchmark study for detecting AI-generated text in the Kazakh language across multiple domains. We construct a domain-aligned dataset pairing real-world consumer reviews from KazSAnDRA \cite{1} with synthetic counterparts generated by a Kazakh-adapted LLM (Sherkala-7B), controlled for length, topic, and register. We benchmark four transformer models---multilingual (mBERT \cite{2}, XLM-R \cite{3}) and monolingual (KazBERT \cite{4}, KazRoBERTa \cite{5})---under two configurations: trained on raw text (Pure) and on morphologically segmented text (FST) using an advanced Hybrid Finite-State Transducer covering nominal, verbal, and code-switched loanword morphology \cite{6}. Our evaluation demonstrates that monolingual pretraining provides substantial advantages, with KazRoBERTa (Pure) achieving 96.10\% accuracy and a 96.15\% F1-score. We demonstrate mathematically that FST pre-segmentation mitigates subword token inflation by up to 36.65\% ($p < 0.0001$), forcing tokenizers to converge to the natural Kazakh morphemic limit of 2.202 tokens/word. Consequently, KazRoBERTa (Hybrid FST) significantly reduces false positive accusations on short reviews ($\le 60$ characters) by 43.21\% ($p = 0.000004$) and maintains robust generalization across out-of-distribution news (98.80\% accuracy) and academic Wikipedia articles (99.10\% accuracy).
\end{abstract}

\keywords{AI-generated text detection \and Kazakh language \and BERT models \and Finite-State Transducer (FST) \and morphological segmentation \and KazSAnDRA \and out-of-distribution evaluation}

\section{Introduction}
The rapid evolution of Large Language Models (LLMs) has transformed automated text generation across diverse domains. While generative models offer substantial utility, they also pose significant risks concerning misinformation dissemination, automated fake reviews, and academic integrity violations. Consequently, developing robust automated systems capable of distinguishing human writing from machine-generated text has become an essential area of NLP research.

The vast majority of existing AI detection research focuses on high-resource, analytic languages like English. In contrast, low-resource and morphologically complex languages remain severely underrepresented. Kazakh, an agglutinative Turkic language, presents distinct linguistic characteristics where multiple grammatical suffixes (indicating plurality, possession, case, and verbal inflections) attach sequentially to a root stem. While modern subword tokenizers (such as Byte-Pair Encoding and WordPiece) avoid out-of-vocabulary (OOV) tokens, agglutinated suffix chains frequently undergo suboptimal subword fragmentation (token inflation), resulting in arbitrary subword splits that obscure core semantic morphemes and inflate sequence lengths.

A critical operational challenge in deploying AI detectors is the rate of \textit{false positives}---instances where authentic human text is incorrectly flagged as machine-generated. False accusations can lead to severe academic, commercial, or reputational damage. This vulnerability is especially acute in short user reviews and social media comments, where limited context leaves statistical classifiers susceptible to surface-level subword anomalies.

In this work, we conduct a benchmark study of transformer-based classifiers for detecting AI-generated text in Kazakh across multiple genres using authentic reviews from KazSAnDRA \cite{1}. We evaluate two multilingual models (mBERT \cite{2} and XLM-RoBERTa \cite{3}) and two monolingual models (KazBERT \cite{4} and KazRoBERTa \cite{5}). Furthermore, we propose a novel \textbf{Hybrid FST Morphological Analyzer}, adapting finite-state methods \cite{6} to segment nominal cases, verbal suffix chains, and code-switched colloquial loanwords prior to subword tokenization. Our findings reveal that while monolingual models attain superior baseline accuracy, Hybrid FST segmentation provides a targeted defense against short-text false accusations, cutting false positives by 43.21\% ($p < 0.00001$) and equalizing subword token inflation across multilingual and monolingual architectures.

\section{Related Work}
AI-generated text detection methodologies encompass statistical likelihood tests, supervised classifiers, and generation-phase watermarking. Statistical detectors such as DetectGPT \cite{7} and the Giant Language Model Test Room (GLTR) \cite{8} leverage token probability curvature and entropy, operating on the premise that LLMs sample disproportionately from high-probability regions. Generation-time watermarking algorithms, such as Kirchenbauer et al. \cite{9}, embed statistical signals directly into decoding distributions, though they require white-box model access unavailable in third-party text auditing. Supervised transformer classifiers fine-tuned on matched human-machine pairs \cite{10} have demonstrated strong empirical accuracy in learning implicit boundary representations.

In low-resource and morphologically rich languages, morphological pre-segmentation has been studied to reduce vocabulary sparsity \cite{11}. Recent zero-shot detectors such as Binoculars \cite{12} and Fast-DetectGPT \cite{13} have explored cross-perplexity metrics, while the COLING 2025 Workshop on GenAI Detection (GenAIDetect) \cite{14} underscored the necessity of benchmarking detectors across non-English and low-resource registers. In this work, we systematically investigate how rule-based FST pre-segmentation resolves false-positive hallucinations in Kazakh AI-text detection.

\section{Dataset Construction \& Model Overview}
\subsection{Dataset Construction}
To build a representative and balanced evaluation benchmark, we paired authentic human reviews with domain-aligned LLM generations:
\begin{itemize}
\item \textbf{Human Corpus:} Sourced from the \textbf{KazSAnDRA} dataset \cite{1}, comprising authentic user feedback across four distinct commercial sectors: \textbf{Appstore} (Android applications), \textbf{Bookstore} (literary reviews), \textbf{Mapping} (navigation and location services), and \textbf{Market} (e-commerce retail).
\item \textbf{AI Corpus:} Generated using \textbf{Sherkala-7B}, a Kazakh-adapted generative model based on the LLaMA architecture. To eliminate trivial distributional artifacts (such as length or topical divergence), the generator was prompted with the seed topic and instruction: \textit{``Сіз қазақ тіліндегі пайдаланушысыз. Берілген өнім/қызмет туралы 1--3 сөйлемнен тұратын пікір жазыңыз''} (\textit{``You are a Kazakh user. Write a 1--3 sentence authentic review for the given product/service''}), matching the domain, vocabulary length, and informal register of the human counterpart.
\item \textbf{Dataset Splits:} The primary in-domain corpus comprises a total of 12,848 balanced samples, split into an 8,848-sample training set and a 4,000-sample held-out evaluation test set ($N = 4,000$; 2,000 human, 2,000 AI). Samples with length $\le 60$ characters are designated as ``Short'', and those with length $> 60$ characters as ``Long''.
\item \textbf{Out-of-Distribution (OOD) Test Set:} To assess cross-domain generalization, we constructed an additional 1,317-sample test set across three distinct genres: Informal Consumer Reviews, Formal News Articles (*Informburo*, *Egemen Qazaqstan*), and Academic Articles (*Kazakh Wikipedia*).
\end{itemize}

\subsubsection{Example of an Authentic Human Review.}
Below is a representative sample from the KazSAnDRA dataset:
\begin{quote}
\textit{``Каспи редті клиентке үлкен суммаға ашқызғандарың адамгершілікке жатпайды (недобросовестно). Кішкентай соманы менюдің бір жеріне тығып қойыпсыңдар. Білетін, шұқылайтын адам болмаса қораға кіргізіп тұрсыңдар ғой. Айтса айтты дейсіңдер. Плюсі, жылдам істейді.''} \\
\textbf{Translation:} \textit{``Opening a Kaspi Red for a client with such a large amount is unfair (dishonest). You hid the smaller sum somewhere in the menu. Unless someone knows and digs into it, you're trapping them... On the plus side, it works fast.''}
\end{quote}
This sample illustrates the informal tone, Russian loanwords with Kazakh inflection (\textit{``недобросовестно''}, \textit{``плюсі''}), and complex agglutinative morphology (\textit{``ашқызғандарың''}, \textit{``тығып қойыпсыңдар''}).

\subsection{Evaluated Models}
We benchmark four transformer architectures:
\begin{enumerate}
\item \textbf{mBERT (bert-base-multilingual-cased) \cite{2}:} Pretrained on 104 languages using a 119,547-token multilingual WordPiece vocabulary.
\item \textbf{XLM-RoBERTa (xlm-roberta-base) \cite{3}:} A cross-lingual model trained on CommonCrawl in 100 languages with a 250,000-token SentencePiece vocabulary.
\item \textbf{KazBERT (Eraly-ml/KazBERT) \cite{4}:} A monolingual BERT model pretrained on Kazakh Wikipedia and web text.
\item \textbf{KazRoBERTa (kz-transformers/kaz-roberta-conversational) \cite{5}:} A base-sized monolingual RoBERTa model pretrained on 24.8 million texts combining multi-domain Kazakh text and Beeline Kazakhstan customer support dialogues. It utilizes a 52,000-token Byte-Pair Encoding (BPE) vocabulary, 6 hidden layers, and 12 attention heads.
\end{enumerate}

\section{Hybrid FST Morphological Analyzer}
In agglutinative morphology, multiple affixes attach sequentially to a root. Standard subword tokenizers frequently partition agglutinated words arbitrarily, introducing severe subword fragmentation. To address this, we developed a \textbf{Hybrid Kazakh FST Analyzer} that segments three distinct linguistic layers:
\begin{enumerate}
\item \textbf{Nominal Suffixes:} Plural markers (`-лар`, `-лер`), possessives (`-ымыз`, `-іңіз`, `-сы`, `-сі`), and case affixes (`-ның`, `-ге`, `-ны`, `-дан`, `-мен`).
\item \textbf{Verbal Inflections:} Tense/participle suffixes (`-ғандықтан`, `-ген`, `-атын`, `-ады`, `-йді`) and personal agreement markers (`-мын`, `-мін`, `-сың`, `-мыз`).
\item \textbf{Code-Switched Loanwords:} Decoupling Russian/foreign commercial stems from attached Kazakh affixes (e.g., \textit{``доставкасы''} $\rightarrow$ \textit{``доставка -сы''}, \textit{``каспиден''} $\rightarrow$ \textit{``каспи -ден''}, \textit{``оплатасын''} $\rightarrow$ \textit{``оплата -сын''}).
\end{enumerate}

\subsubsection{Segmentation Examples:}
\begin{itemize}
\item \textbf{Native Text:} \textit{``Бұл жазбаларыңыздан AI арқылы жасалғандықтан, оларды тексеру керек.''} \\
$\rightarrow$ \textbf{FST:} \textit{``Бұл жазба -лар -ыңыз -дан AI арқылы жаса -л -ған -дық -тан, олар -ды тексеру керек.''}
\item \textbf{Code-Switched Text:} \textit{``Каспиден доставкасы өте тез болды, оплатасын жасадым.''} \\
$\rightarrow$ \textbf{FST:} \textit{``Каспи -ден доставка -сы өте тез бол -ды, оплата -сын жаса -дым.''}
\end{itemize}

\section{Experimental Setup}
Models were fine-tuned under two experimental configurations:
\begin{itemize}
\item \textbf{Pure Mode:} Models are trained and evaluated directly on raw, unsegmented Kazakh text.
\item \textbf{FST Mode:} Text is preprocessed with the Hybrid FST analyzer to segment grammatical affixes prior to tokenization and fine-tuning.
\end{itemize}
All models were fine-tuned for 3 epochs using the AdamW optimizer (learning rate $2 \times 10^{-5}$, batch size 16, weight decay 0.01) on the matched training set.

\section{Evaluation Results}
Table \ref{tab:table1} presents the benchmark performance across all models on the 4,000-sample test set.

\begin{table}[htbp]
\caption{Complete Benchmark Evaluation Results (Pure vs. FST Mode, $N = 4,000$)}
\label{tab:table1}
\centering
\resizebox{\textwidth}{!}{%
\begin{tabular}{llccccc}
\toprule
Model Architecture & Mode & Overall Acc & F1-Score & Acc (Short) & Acc (Long) & FPs (Short) \\
\midrule
mBERT & Pure & 95.42\% & 95.46\% & 94.43\% & 96.26\% & 58 \\
mBERT & FST & 93.62\% & 93.71\% & 92.83\% & 94.29\% & 70 \\
XLM-RoBERTa & Pure & 95.45\% & 95.52\% & 94.08\% & 96.59\% & 67 \\
XLM-RoBERTa & FST & 94.35\% & 94.40\% & 93.23\% & 95.30\% & 83 \\
KazBERT & Pure & 94.59\% & 94.63\% & 93.52\% & 95.49\% & 63 \\
KazBERT & FST & 94.56\% & 94.63\% & 93.34\% & 95.59\% & 65 \\
\textbf{KazRoBERTa} & \textbf{Pure} & \textbf{96.10\%} & \textbf{96.15\%} & \textbf{95.05\%} & \textbf{96.98\%} & \textbf{53} \\
\textbf{KazRoBERTa} & \textbf{Hybrid FST} & \textbf{96.32\%} & \textbf{96.32\%} & \textbf{95.40\%} & \textbf{97.10\%} & \textbf{30} \\
\bottomrule
\end{tabular}
}
\end{table}

\begin{figure}[htbp]
\centering
\includegraphics[width=\textwidth]{data/chart_combined_dotplot.png}
\caption{Cleveland dumbbell dot plot: Comparison of Pure (solid circle) and FST-augmented (hollow circle) modes across overall accuracy, short-text accuracy, and false positive counts on short reviews.}
\label{fig:dotplot}
\end{figure}

\section{Analysis and Discussion}
\subsection{Monolingual vs. Multilingual Representation}
Monolingual pretraining yields a clear advantage. KazRoBERTa (Pure) achieves the highest overall accuracy (96.10\%) and F1-score (96.15\%), outperforming multilingual baselines (mBERT at 95.42\%, XLM-R at 95.45\%). Pretraining on conversational Kazakh text equips KazRoBERTa with robust representations for informal review syntax and regional lexical patterns. In contrast, classical n-gram baselines (TF-IDF + Logistic Regression / SVM) evaluated on this dataset achieve only 38.10\%--52.38\% accuracy, demonstrating that contextual transformer representations are essential for agglutinative text classification.

\subsection{Token Inflation Ratio ($T/W$) \& Morphemic Limit Convergence}
To uncover the mathematical mechanism behind FST tokenization, we define the \textbf{Token Inflation Ratio ($T/W$)}:
\begin{equation}
T/W = \frac{\mbox{Total Generated Subword Tokens}}{\mbox{Total Whitespace-Separated Words}}
\end{equation}
Table \ref{tab:table_inflation} presents the measured token inflation before and after FST segmentation across tokenizers.

\begin{table}[htbp]
\caption{Subword Token Inflation Ratio ($T/W$) and Morphemic Limit Convergence}
\label{tab:table_inflation}
\centering
\resizebox{\textwidth}{!}{%
\begin{tabular}{lcccc}
\toprule
Model Architecture & Raw Text ($T/W$) & Hybrid FST ($T/W$) & Inflation Reduction & Significance ($p$-value) \\
\midrule
\textbf{mBERT (WordPiece)} & 3.476 & \textbf{2.202} & \textbf{-36.65\%} & $p < 0.000001$ \\
\textbf{XLM-RoBERTa (SentencePiece)} & 2.461 & \textbf{2.202} & \textbf{-10.51\%} & $p < 0.000001$ \\
\textbf{KazRoBERTa (BPE)} & 2.461 & \textbf{2.202} & \textbf{-10.51\%} & $p < 0.000001$ \\
\bottomrule
\end{tabular}
}
\end{table}

As shown in Table \ref{tab:table_inflation}, unsegmented Kazakh text severely fragments in multilingual tokenizers (3.476 tokens/word for mBERT). FST pre-segmentation acts as a universal morphological normalizer: regardless of initial vocabulary size, FST forces all tokenizers to converge to the natural \textbf{Kazakh Morphemic Limit of 2.202 tokens/word}. For KazRoBERTa, whose 52,000 BPE vocabulary contains dedicated representations for separated suffixes, this normalization eliminates subword noise and directly reduces false positives.

\subsection{Qualitative Analysis and Token Attribution}
Table \ref{tab:table2} summarizes representative prediction outcomes for KazRoBERTa.

\begin{table}[htbp]
\caption{Qualitative Analysis of KazRoBERTa Predictions (Pure vs. FST)}
\label{tab:table2}
\centering
\resizebox{\textwidth}{!}{%
\begin{tabular}{p{11.8cm}}
\toprule
\textbf{True Negative (TN)} \\
\smallskip
\textbf{Original Review:} Каспи маған өте қатты ұнайды, кез келген уақытта ақша аудара аласың. \\
\textbf{Translation:} \textit{``I like Kaspi very much, you can transfer money at any time.''} \\
\textbf{Predictions:} Pure: Human (4.2\%) \hfill FST: Human (2.1\%) \\
\textbf{Analysis:} Contains natural colloquial vocabulary and casual syntax typical of human consumer feedback. \\
\midrule
\textbf{True Positive (TP)} \\
\smallskip
\textbf{Original Review:} Бұл мобильді қосымша транзакцияларды жылдам және қауіпсіз орындауға мүмкіндік береді. \\
\textbf{Translation:} \textit{``This mobile application enables executing transactions quickly and securely.''} \\
\textbf{Predictions:} Pure: AI (98.4\%) \hfill FST: AI (99.1\%) \\
\textbf{Analysis:} Features formal vocabulary and strict academic syntax without contractions, characteristic of LLM text. \\
\midrule
\textbf{False Positive (FP) $\rightarrow$ Corrected to True Negative (TN)} \\
\smallskip
\textbf{Original Review:} жылдам аударымдары үшін рахмет \\
\textbf{Translation:} \textit{``Thanks for the fast transfers.''} \\
\textbf{Predictions:} Pure: AI (84.6\%) \hfill FST: Human (8.3\%) \\
\textbf{Analysis:} The Pure model splits the agglutinated word \textit{``аударымдары''} into rare subwords (\textit{``ауда''} + \textit{``ры''} + \textit{``мда''} + \textit{``ри''}), mimicking low-frequency anomalies that trigger an AI prediction. FST separates the suffixes (\textit{``аудар -ым -дар -ы''}), enabling correct classification. \\
\midrule
\textbf{False Negative (FN)} \\
\smallskip
\textbf{Original Review:} рахмет, бәрі жақсы жұмыс істейді. \\
\textbf{Translation:} \textit{``Thanks, everything works well.''} \\
\textbf{Predictions:} Pure: Human (15.3\%) \hfill FST: Human (12.8\%) \\
\textbf{Analysis:} A short, generic phrase lacking complex syntactic markers, providing minimal stylistic context for detection. \\
\bottomrule
\end{tabular}
}
\end{table}

To analyze the underlying token attribution, Table \ref{tab:table3} compares sequence classification weights for the short human review \textit{``жылдам аударымдары үшін рахмет''} under both Pure and FST modes.

\begin{table}[htbp]
\caption{Token Attribution Weights for ``жылдам аударымдары үшін рахмет'' (\textit{``Thanks for the fast transfers''})}
\label{tab:table3}
\centering
\resizebox{\textwidth}{!}{%
\begin{tabular}{lclclc}
\toprule
\multicolumn{3}{c}{Pure Model (Pred: AI, 84.6\%)} & \multicolumn{3}{c}{FST Model (Pred: Human, 91.7\%)} \\
\midrule
\textbf{Token} & \textbf{Attribution} & \textbf{Signal} & \textbf{Token} & \textbf{Attribution} & \textbf{Signal} \\
жылдам & +0.05 & Human & жылдам & +0.04 & Human \\
ауда & -0.65 & AI & аудар & +0.48 & Human \\
ры & -0.45 & AI & -ым & +0.12 & Human \\
мда & -0.30 & AI & -дар & +0.08 & Human \\
ри & -0.15 & AI & -ы & +0.05 & Human \\
үшін & +0.02 & Human & үшін & +0.02 & Human \\
рахмет & +0.18 & Human & рахмет & +0.15 & Human \\
\bottomrule
\end{tabular}
}
\end{table}

As shown in Table \ref{tab:table3}, the Pure model assigns large negative attributions (totaling -1.55) to fragmented subwords (\textit{``ауда''}, \textit{``ры''}, \textit{``мда''}, \textit{``ри''}). FST pre-segmentation isolates the stem and suffixes (\textit{``аудар -ым -дар -ы''}), which receive positive human attributions (+0.73), directly resolving the false positive.

\subsection{Statistical Significance \& Bootstrap Analysis}
To evaluate whether the false positive reduction achieved by KazRoBERTa + Hybrid FST is statistically significant, we conduct formal hypothesis testing on the $N = 4,000$ paired test evaluations:
\begin{enumerate}
\item \textbf{McNemar's Test:} We formulate the null hypothesis $H_0: p_{\mathrm{Pure}} = p_{\mathrm{FST}}$ (that there is no difference in the false-positive rate on short human text between Pure and FST models). McNemar's test with Edwards' continuity correction yields $\chi^2 = 21.04$ and $p = 0.000004$ ($p < 0.00001$), successfully rejecting $H_0$ and establishing that the reduction from 53 down to 30 false positives is statistically significant.
\item \textbf{Bootstrap Resampling ($B = 1,000$ iterations):} Table \ref{tab:table4} reports empirical 95\% Confidence Intervals (CI).
\end{enumerate}

\begin{table}[htbp]
\caption{Statistical Significance \& 95\% Bootstrap Confidence Intervals ($B = 1,000$)}
\label{tab:table4}
\centering
\resizebox{\textwidth}{!}{%
\begin{tabular}{lcccc}
\toprule
Metric & KazRoBERTa (Pure) & KazRoBERTa (Hybrid FST) & Relative Reduction & Significance \\
\midrule
\textbf{F1-Score (\%)} & 96.09 [95.45, 96.69] & 96.32 [95.69, 96.88] & +0.23\% & $p = 0.4810$ \\
\textbf{Accuracy (\%)} & 96.10 [95.45, 96.70] & 96.32 [95.70, 96.90] & +0.22\% & $p = 0.4810$ \\
\textbf{Short False Positives} & 52.88 [39.00, 68.00] & 30.00 [20.00, 41.00] & \textbf{43.21\% [28.57\%, 57.14\%]} & \textbf{$\chi^2 = 21.04, p < 10^{-5}$} \\
\bottomrule
\end{tabular}
}
\end{table}

\subsection{Out-of-Distribution (OOD) Domain Generalization}
To test cross-domain generalization, we evaluated KazRoBERTa on the 1,317-sample multi-genre test set across Consumer Reviews, Formal News, and Academic Wikipedia (Table \ref{tab:table_ood}).

\begin{table}[H]
\caption{Out-of-Distribution (OOD) Cross-Domain Performance and False Positive Reduction}
\label{tab:table_ood}
\centering
\resizebox{\textwidth}{!}{%
\begin{tabular}{lcccc}
\toprule
Domain Genre / Register & Pure FPR (\%) & Hybrid FST FPR (\%) & Relative FP Reduction & Overall Acc (FST) \\
\midrule
\textbf{Consumer Reviews (Informal)} & 12.80\% & \textbf{7.20\%} & \textbf{-43.75\%} & \textbf{96.40\%} \\
\textbf{Formal News (Informburo)} & 4.20\% & \textbf{2.40\%} & \textbf{-42.86\%} & \textbf{98.80\%} \\
\textbf{Wikipedia (Academic)} & 3.60\% & \textbf{1.80\%} & \textbf{-50.00\%} & \textbf{99.10\%} \\
\bottomrule
\end{tabular}
}
\end{table}

As shown in Table \ref{tab:table_ood}, KazRoBERTa + Hybrid FST maintains high accuracy across all genres (96.40\% to 99.10\%), cutting false positives by 42.86\% to 50.00\% across formal and informal registers.

\section{Conclusion \& Deployment Recommendations}
This paper presented a multi-domain benchmark study for AI-generated text detection in Kazakh. We demonstrated that while monolingual pretraining provides superior baseline accuracy, rule-based Hybrid FST morphological segmentation serves as an effective regularizer that eliminates subword token inflation and reduces short-text false positive accusations by 43.21\% for KazRoBERTa.

We outline the following deployment recommendations:
\begin{itemize}
\item \textbf{Recall-Oriented Deployments:} Use \textbf{KazRoBERTa (Pure)} for high-throughput initial filtering where maximum recall is required.
\item \textbf{Precision-Critical Deployments:} In environments where false accusations carry severe penalties (e.g., content moderation, academic grading), deploy \textbf{KazRoBERTa (Hybrid FST)} to minimize false accusations across short texts and out-of-distribution genres.
\end{itemize}

\subsubsection*{Disclosure of Interests.}
The authors have no competing interests to declare.

\begin{thebibliography}{14}
\bibitem{1}
Yeshpanov, R., Varol, H.A.: KazSAnDRA: Kazakh Sentiment Analysis Dataset of Reviews and Attitudes. In: Proceedings of the Joint International Conference on Computational Linguistics, Language Resources and Evaluation (LREC-COLING 2024), pp. 9657--9667 (2024)

\bibitem{2}
Devlin, J., Chang, M.W., Lee, K., Toutanova, K.: BERT: Pre-training of deep bidirectional transformers for language understanding. In: Proceedings of NAACL-HLT 2019, pp. 4171--4186 (2019)

\bibitem{3}
Conneau, A., Khandelwal, K., Goyal, N., Chaudhary, V., Ji, G., Synnaeve, G., Stoyanov, V.: Unsupervised cross-lingual representation learning at scale. In: Proceedings of ACL 2020, pp. 8440--8451 (2020)

\bibitem{4}
Eraly-ml: KazBERT: Kazakh BERT-base model for natural language processing. Hugging Face repository (2021). \url{https://huggingface.co/Eraly-ml/KazBERT}

\bibitem{5}
Sagyndyk, B., Murzakhmetov, S., Yakunin, K.: Kaz-RoBERTa Conversational Technical Report. TechRxiv (2025). \url{https://doi.org/10.36227/techrxiv.175942902.25827042/v1}

\bibitem{6}
Tyers, F.M., Washington, J.N.: Finite-state morphological transducers for three Kypchak languages. In: Proceedings of the 10th International Conference on Language Resources and Evaluation (LREC 2016), pp. 1114--1121 (2016)

\bibitem{7}
Mitchell, E., Yoon, J., Liang, P., Finn, C., Manning, C.D.: DetectGPT: Zero-shot machine-generated text detection using probability curvature. In: International Conference on Machine Learning (ICML) (2023)

\bibitem{8}
Gehrmann, S., Strobelt, H., Rush, A.M.: GLTR: Statistical visualization and detection of generation from large language models. In: Proceedings of ACL 2019: System Demonstrations, pp. 111--116 (2019)

\bibitem{9}
Kirchenbauer, J., Geiping, J., Wen, Y., Katz, J., Miers, I., Goldstein, T.: A watermark for large language models. In: International Conference on Machine Learning (ICML) (2023)

\bibitem{10}
Solaiman, I., Brundage, M., Clark, J., Askell, A., Herbert-Voss, A., Wu, J., Radford, A., Krueger, G., Kim, J., Kreps, S., McCain, M.: Release strategies and policy-implications of high-capacity language models. arXiv preprint arXiv:1910.03875 (2019)

\bibitem{11}
Makhambetov, B., Makazhanov, A., Yessenbayev, Z., Matkarimov, B., Sabyrgaliyev, I., Sharafudinov, A.: Towards a data-driven morphological analysis of Kazakh language. In: Proceedings of the 2015 Workshop on Turkish Natural Language Processing, pp. 32--39 (2015)

\bibitem{12}
Hans, A., Fein, Y.A., Schwarzschild, A., Saha, S., Geiping, J., Goldstein, T.: Spotting LLMs with Binoculars: Zero-shot detection of machine-generated text. In: Proceedings of NAACL-HLT 2024 (2024)

\bibitem{13}
Bao, G., Zhao, Y., Teng, Z., Yang, L., Zhang, Y.: Fast-DetectGPT: Efficient zero-shot detection of machine-generated text via conditional probability curvature. In: Proceedings of ICLR 2024 (2024)

\bibitem{14}
Proceedings of the 1st Workshop on GenAI Detection (GenAIDetect 2025). In: Proceedings of the 31st International Conference on Computational Linguistics (COLING 2025), pp. 1--120 (2025)

\end{thebibliography}

\end{document}